% This is samplepaper.tex, a sample chapter demonstrating the
% LLNCS macro package for Springer Computer Science proceedings;
% Version 2.21 of 2022/01/12
%
\documentclass[runningheads]{llncs}
%
\usepackage[T1]{fontenc}
% T1 fonts will be used to generate the final print and online PDFs,
% so please use T1 fonts in your manuscript whenever possible.
% Other font encondings may result in incorrect characters.
%
\usepackage[utf8]{inputenc}
\usepackage{url}
\usepackage{graphicx}
\usepackage{hyperref}
\usepackage{listings}
\usepackage{amsmath}
\usepackage{amssymb}
\usepackage[inline]{enumitem}
\usepackage{makecell}
\usepackage{tabularx}
% Used for displaying a sample figure. If possible, figure files should
% be included in EPS format.
%
% If you use the hyperref package, please uncomment the following two lines
% to display URLs in blue roman font according to Springer's eBook style:
%\usepackage{color}
%\renewcommand\UrlFont{\color{blue}\rmfamily}
%\urlstyle{rm}
%

\begin{document}
    \setlist[enumerate]{noitemsep, topsep=0pt}
%
%    \title{{\fontsize{14}{16}\selectfont Client-Driven Offline-First RDF 1.2 using OR-Sets}}

    \title{Client-Driven Offline-First RDF 1.2 using OR-Sets}
%
%\titlerunning{Abbreviated paper title}
% If the paper title is too long for the running head, you can set
% an abbreviated paper title here
%
    \author{Jitse De Smet\inst{1}\orcidID{0009-0002-6513-5013} \and
    Ruben Taelman\inst{1}\orcidID{0000-0001-5118-256X}}
%
    \authorrunning{J. De Smet and R. Taelman}
% First names are abbreviated in the running head.
% If there are more than two authors, 'et al.' is used.
%
    \institute{IDLab, Department of Electronics and Information Systems, Ghent University – imec.
    \email{\{firstname\}.\{lastname\}@ugent.be}}
%
    \maketitle              % typeset the header of the contribution
%
    \begin{abstract}
        As autonomous agents increasingly write to shared RDF knowledge graphs concurrently,
        especially in the context of decentralized data ecosystems, principled conflict resolution becomes critical.
        Existing RDF CRDT approaches either require dedicated infrastructure or introduce
        consistency boundary problems when storing bookkeeping across multiple resources.
        We instantiate the state-based OR-Set over RDF 1.2 triple terms -
        unlike prior work, retaining all bookkeeping within a single dataset without additional infrastructure -
        and demonstrate a proof-of-concept datastore wrapper to use in query engines,
        enabling agents to query and update a conflict-free knowledge graph via standard SPARQL,
        unaware of the underlying merge machinery.
        Future work will investigate constraint-aware merge semantics for RDF CRDTs.


        \keywords{Conflict-free Replicated Data Types \and
                  OR-Set \and
                  RDF 1.2 \and
                  Decentralized Web \and
                  Offline-First \and
                  SPARQL}

        \textbf{npm package}: \href{https://www.npmjs.com/package/orset-rdf-store}{https://www.npmjs.com/package/orset-rdf-store} \\
        \textbf{poster}: \href{https://orset-rdf.poster.jitsedesmet.be/}{https://orset-rdf.poster.jitsedesmet.be/}
    \end{abstract}
%
%
%


    \section{Introduction}
    \label{introduction}

    Knowledge graphs expressed in RDF are a central pillar of neuro-symbolic AI:
    techniques such as GraphRAG~\cite{graphrag}, GRASP~\cite{grasp}, and the Model Context Protocol (MCP)~\cite{mcp}
    make it increasingly common for multiple agents to read from and write to the same RDF knowledge graph simultaneously -
    one agent enriches the graph while another reasons over it, and a third corrects errors it discovers.

    In decentralized data ecosystems - such as Solid~\cite{solid}, IDSA~\cite{idsa}, and Gaia-X~\cite{gaia-x} -
    each participant governs their own data without relying on a central authority.
    A critical challenge is concurrent offline editing:
    when two agents independently modify the same resource while offline, their updates collide once connection is restored.
    Asking an agent to resolve such conflicts manually is not always desirable - an agent that has been offline may struggle to grasp how much has changed.
    Conflict-free Replicated Data Types (CRDTs)~\cite{CRDTs} offer a principled solution,
    guaranteeing that any two replicas converge to the same state through well-defined merge algorithms, with no coordination required.

    Solutions have been proposed for RDF-based CRDTs~\cite{rdf-crdt} and even CRDTs in Solid~\cite{solid-crdt-yjs},
    but the introduction of RDF 1.2~\cite{rdf-1-2} opens new possibilities.
    RDF 1.2's triple terms - where the object of a triple can itself be a triple - enable efficient,
    first-class modelling of state-based CRDTs directly within RDF datasets, without any external bookkeeping.

    In this paper, we present a\textbf{state-based add-wins set CRDT} modelled entirely in RDF 1.2.
    All bookkeeping stays within the same dataset,
    enabling atomic updates on the resource without server intervention,
    and allowing both \textbf{CRDT-aware and CRDT-unaware clients} to coexist without friction.
    By plugging a CRDT-aware datastore into a query engine, conflict resolution becomes fully transparent to both agents and query engines -
    as demonstrated by our \href{https://www.npmjs.com/package/orset-rdf-store}{PoC implementation}.
    Agents can then query and update a conflict-free knowledge graph through standard SPARQL interfaces,
    without any awareness of the underlying CRDT machinery.

    \section{Related Work}
    \label{sota}

    \textbf{Decentralized Data Ecosystems and Interface Heterogeneity}
    Decentralized data ecosystems consist of self-governed data stores,
    each individually positioned in the continuous \textit{Consistency Availability Partition-tolerance} (CAP) space and exposed through heterogeneous interfaces~\cite{phd-symposium} -
    their only shared infrastructure being the Web itself.
    Our CRDT is designed to work across such ecosystems by depending only on what every participant already has:
    the ability to store and exchange RDF data over the Web, plus \textbf{ETag support} to avoid
    \href{https://developer.mozilla.org/en-US/docs/Web/HTTP/Reference/Headers/ETag}{mid-air collisions}.


    \textbf{CRDTs}
    Existing RDF CRDT approaches require more infrastructure than the Web alone.
    The m-ld project~\cite{m-ld} provides a JavaScript CRDT engine for RDF with a similar eventual-consistency goal,
    but requires a dedicated m-ld domain server.
    Gruss et al.~\cite{solid-crdt-yjs} store the latest RDF state alongside a binary CRDT representation from a user-chosen library
    (e.g. \href{https://yjs.dev/}{Yjs} or \href{https://automerge.org/}{Automerge}),
    plus a hypermedia description of supported operations.
    While this lets non-CRDT-aware clients consume the data, it relies on a separate binary blob outside the RDF dataset,
    introducing a consistency boundary problem where a standard HTTP server cannot atomically update multiple resources.
    Braid-HTTP~\cite{braid} proposes adding merge-type headers to HTTP,
    which is complementary to our approach and could replace our custom patch encoding in a future iteration.

    Since we abstract at the RDF dataset level - where only set-based operations exist - we limit ourselves to set CRDTs.
    The SU-Set~\cite{rdf-crdt} extends the OR-Set~\cite{CRDTs} for RDF using an operation-based approach,
    whereas we adopt the state-based OR-Set formalization of Bieniusa et al.~\cite{optimized-set-crdt},
    which aligns naturally with HTTP's Representational State Transfer (REST) semantics -
    exchanging full dataset state rather than individual operations.
    Where Bieniusa et al. optimize tombstone removal using causal delivery and a known replica set,
    we instead exploit NTP time drift bounds - an approach better suited to open decentralized ecosystems where replica sets are unknown.


    \section{State-based Add-Wins CRDT-RDF}
    \label{core}

    The OR-Set~\cite{optimized-set-crdt} is a replicated set where each element carries a set of unique add-tags;
    removing an element moves its tags to a tombstone set rather than deleting them outright,
    ensuring that a concurrent add always wins over a concurrent remove.
    We instantiate this over RDF 1.2, setting elements as RDF 1.2 triples and unique tags as \href{https://www.rfc-editor.org/rfc/rfc9562.html}{UUDv4} identifiers.
    RDF 1.2 triple terms allow a triple to appear in the object position of another triple,
    enabling statements to be made about triples directly.
    We use this to link each tracked triple, its add-tags, and its tombstones through an identifying node,
    as shown in \autoref{fig-transformer}.
    A triple is considered present when it has at least one add-tag not in its tombstone set.
    On merge, multiple identifying nodes for the same triple are consolidated into one.

    \begin{figure}[tbp]
    \centering
        \begin{lstlisting}[language=JavaScript,basicstyle=\fontsize{6}{7.6}\ttfamily,breaklines=true]
<> a crdt:container .
[] crdt:tags <<( <> a crdt:container )>> ;
   crdt:add "be2f95dd-8ca9-416c-b1d0-81dc45ba54c8"^^crdt:uuid .
:me a :human .
[] crdt:tags <<( :me a :human )>> ;
   crdt:add "96d482e5-3ce7-4f24-a21b-a9d0506ff5b0"^^crdt:uuid ;
   crdt:remove "77c01067-1594-475e-8c64-76f9c4ec4402"^^crdt:uuid .
[] crdt:tags <<( :me a :man )>> ;
   crdt:remove
"216ac011-c2ba-4ff0-825c-7f9cf1efa4ff--2025-03-13T14:00:00Z"^^crdt:stamp-uuid .
        \end{lstlisting}
        \caption{RDF 1.2 representation of a state-based OR-Set with two triples present:
                     `<> a crdt:container` and `:me a :human`.
                     The latter has one add-tag not in its tombstone set, so it remains present.
                     Since `:me a :man` has no add-tags, it is not present.}
        \label{fig-transformer}
        \vspace{-12pt}
    \end{figure}

    \subsubsection{NTP Tombstone Optimization}

    Without garbage collection, tombstones accumulate indefinitely -
    a fundamental challenge for all state-based CRDTs.
    Bieniusa et al.'s solution requires causal delivery and a known replica set,
    assumptions that do not hold in open decentralized ecosystems.
    We instead exploit  \href{https://www.rfc-editor.org/rfc/rfc5905}{NTP} synchronization:
    since NTP enters panic mode after a drift exceeding \(1000\text{s}\),
    any two synchronized clients differ by at most \(2000\text{s}\).
    By stamping each tag with its creation or tombstone time using `crdt:stamp-uuid` -
    formatted as \texttt{\{uuid\}\{xsd:dateTime\}} -
    a CRDT dataset can declare a synchronization interval \(d\), after which stale metadata is pruned:
    1. an add-tag may be dropped if a newer one exists and all replicas have seen it;
    2. a tombstone may be dropped once all replicas have seen it; and
    3. the identifying blank node may be dropped if no other triple references it.
    A tag is universally seen once its timestamp is older than \(d + 4000\text{s}\) -
    a bound derived from the worst-case NTP drift between any two synchronized replicas.


    \section{Conclusion}
    \label{conclusion}

    We presented a state-based add-wins set CRDT modelled entirely in RDF 1.2,
    requiring no infrastructure beyond RDF-over-HTTP with ETag support - something most modern Web servers already provide.
    All bookkeeping stays within the dataset itself, eliminating the consistency boundary problem.
    An NTP-based tombstone pruning mechanism keeps the CRDT manageable without requiring a known replica set.
    By abstracting at the quad-store level, the CRDT can be plugged into existing query engines,
    making conflict resolution fully transparent to both agents and SPARQL clients alike.
    A current limitation is the lack of constraint-aware merge:
    two individually valid states can merge into a state that violates domain constraints,
    a problem we plan to address in future work.

    \begin{credits}
        \subsubsection{\ackname}
        Jitse De Smet is a predoctoral fellow of the Research Foundation -- Flanders (FWO) (1SB8525N).
        The described research activities were supported by SolidLab Vlaanderen (Flemish Government, EWI and RRF project VV023/10).
        Ruben Taelman is a postdoctoral fellow of the Research Foundation -- Flanders (FWO) (1202124N).
    \end{credits}

%
% ---- Bibliography ----
%
% BibTeX users should specify bibliography style 'splncs04'.
% References will then be sorted and formatted in the correct style.
%
    \bibliographystyle{splncs04}
    \bibliography{references}
\end{document}
